% Transcription — fonds Alexandre Grothendieck, shelfmark 43, batch 10 (pages 181-200).
% Established from the facsimile archives/batches/43/batch-10.pdf (PDF page k =
% archivists' page 180+k, no cover sheet), cut from a copy of 43.pdf fetched
% through the project's own relay
% (https://grothendieck-relay.onrender.com/source/43.pdf); tiles in
% archives/tiles/43/p181-p200. Per the skill transcribe-grothendieck. The folder
% is 207 pages in eleven batches, transcribed in parallel; batch 9 (pages
% 161-180) existed when this file was finished and its notation was read and
% followed (G_m, G_a, alpha_p, mu_p, Z/lZ, the double-struck V for the Witt
% covector group, double-struck W for Witt vectors, \underline for his
% underlined Hom).
%
% Pass: Opus 5.5 (claude-opus-5-5), 2026-09-26 — first pass, unchecked against the pages by a human.
%
% Dating: the folder has its own inventory date, carried verbatim in
% \dating{} with no group sentence (the group « Jacobiennes » (43-44),
% dated [à partir de 1961-à partir de 1967], is added by the skill only for
% undated folders; batch 9 was aligned the same way). No page of this batch carries a date.
%
% Copies. No page of this batch is evidently a photocopy: pages 181-192 are
% blue-black ink on grey-blue paper, 194-195 ink on cream paper, 196-200 ink
% on the back and front of sheets of a typist's printed letterhead (the
% heading is printed at the top left of p. 198 and upside down at the foot of
% pp. 197, 199, 200). If some of 181-192 are copies, nothing in the ink shows
% it at this resolution.
%
% Pages skipped: 193, a blank cream sheet; 197, the blank side of a
% letterhead sheet carrying only the printed heading.
%
% Pages summarised: none. The margins of 186, 188, 190, 191 and 192 carry
% long diagonal notes that are mostly illegible; what reads is given, the
% rest described in one \note per block.
%
% His pagination: none on these twenty leaves; no numbered run. The run of
% batch 9 under the cover « Dualité de Cartier des groupes affines et groupes
% formels » (from p. 170) continues here without a new heading: p. 181 takes
% up the categories C_1, C_2 and A of p. 180. Page 194 opens a new heading of
% his, « Catégories abéliennes », set as the \section.
%
% What the batch is about. Cartier duality and the Dieudonné ring. P. 181:
% pairs (X, Y) of A-modules with H^0_m(X) ≅ D H^0_m(Y) and their morphisms;
% to include G_m and Z one takes Z × A. P. 182: how to define nullity
% (nilpotence?) of F, V intrinsically on ind- and pro-objects; modules over
% W_sigma[F,V]/(FV-p) as endomorphisms of W_infty preserving D(W_infty).
% P. 183 (largely cancelled): 0 -> mu_n -> G_m -> G_m -> 0, Ext^1(G_m,mu_n),
% W_sigma[F,V]/(V^n, FV-p), a boxed scheme of the « elementary » groups in
% opposite orders, « Anneau de Cartier » vs « Anneau de Gabriel-Dieudonné ».
% P. 184: D(W_infty) ⊂ W_infty1, affine groups = algebras of finite type;
% question whether an affine group is pro-(affine algebraic). Pp. 185-188:
% ind-(finite algebraic) groups as coalgebras, Cartier: duality
% pro(linear algebraic groups) ~ gr(ind. alg. finite); G -> G_m versus
% group-like f in A (Delta f = f ⊗ f, eps(f) = 1); the dual D(G) as an
% ind-affine group representing Hom_{k-gr}(G, G_m); D(G_m) = Z,
% D(Z/lZ) = Z/lZ, D(Z/p^nZ) = mu_{p^n}, D(mu_{p^n}) = Z/p^nZ, D(G_a) = ... .
% Pp. 189-190: Pro(C) and the functor Pro(C) -> Hom(C°, Ens), not fully
% faithful (example C = Ens); lim_C X; pro-affine schemes of finite type =
% affine schemes; question: is an affine group scheme a pro(affine group of
% finite type). P. 191: hyperalgebras as inductive limits of sub-hyperalgebras
% of finite type (easy over a field); the « dual » D^(G) whose hyperalgebra is
% the affine algebra of G. P. 192: the elementary groups dual to each other,
% V = D(G_a), a table multiplicative / unipotent / discrete. Pp. 194-195
% (fair copy): Ind C for C abelian with objects of finite length; the socle,
% F_0(X), F_1(X), the filtration Phi_i(X) by Krull dimension. P. 196:
% ind(pro-(finite alg. groups)) = pro(Modules ... ) dual to modules over A;
% a triangle groupes formels raisonnables / pro-(groupes alg. finis) /
% Modules sur A; W_sigma[F][F^{-1}] and F^{-1} w = sigma^{-1}(w) F^{-1}.
% P. 198: W_{n,k} = Ker(F^k on W_n), W̄_n = union, U = lim W̄_i, F, V on U,
% V(wx) = sigma^{-1}(w)V(x), boxed FV = VF = p, F^n U = 0, p^n U = 0.
% P. 199: A = W_infty[[F,V]] with FV = VF = p, Fw = sigma(w)F,
% wV = V sigma(w); left/right multiplication computed on a general element;
% dictionary: finite commutative alg. groups <-> A-modules of finite length,
% pro(finite alg. groups) (noeth.) <-> modules loc. of finite length (art.),
% « groupes formels » (art.) <-> modules of finite type. P. 200: the same
% dictionary « sans isogénie », A_F, and the boxed definition: modules of
% finite type M over W with additive bijective F and additive V, F(wx) =
% F(w)F(x), V(wx) = F^{-1}(w)V(x).
%
% Notation fixed once for the batch: \mathbb{G}_m, \mathbb{G}_a, \mu_p,
% \alpha_p, \mathbb{Z}/\ell\mathbb{Z}; \mathbb{W}_\infty, \mathbb{W}_\sigma
% where he doubles the W (pp. 182-184); plain W where he writes it single
% (pp. 196, 198-200), with a note on p. 198; bold \mathbf{W}, \mathbf{V},
% \mathbf{F} where he inks them heavily (pp. 198, 200); \mathbb{V} for the
% double-struck V of p. 192 (= D(G_a), as in batch 9); D(G) for the Cartier
% dual; \underline{C} for his underlined categories (pp. 189-190, 194-195);
% \Phi for the round Φ of p. 195; \mathcal{U} for the cursive U of p. 198;
% \mathfrak{m} for the idealised m of p. 181.
%
% Readings to check first: nearly all the connective prose of pp. 181-192
% (fast drafting; the formulas carry, the prose often does not); the
% heading letters of p. 181; « nullité » (or « nilpotence ») on p. 182; the
% letter before (F,V) on p. 183; the mid-page of p. 192 and its table; the
% labels « contrav. », « cov. » on pp. 196 and 200; the indices of the
% multiplication computation on p. 199; « sans isogénie » on p. 200.

\documentclass[11pt,a4paper]{article}
% Le préambule commun à toutes les transcriptions du fonds Grothendieck.
%
% Il tient en une idée : **l'incertitude se compose, elle ne se résout pas.**
% Une transcription qui lisse un mot douteux a détruit la seule chose pour
% laquelle on la faisait — un lecteur doit pouvoir savoir, à chaque endroit,
% ce qui était sur le feuillet et ce qui a été ajouté. Les six macros qui
% suivent sont donc l'appareil critique, et non de la décoration.
%
% Les mêmes commandes sont reconnues par `scripts/render.mjs`, qui produit la
% vue de lecture du site. Une macro ajoutée ici doit l'être aussi là-bas,
% faute de quoi le rendu échouera bruyamment — ce qui est voulu : un rendu
% silencieusement faux est pire qu'un rendu absent.

\ProvidesPackage{grothendieck}[2026/08/08 Transcriptions du fonds Grothendieck]

\usepackage{amsmath,amssymb,amsthm}
\usepackage{mathrsfs}
\usepackage{stmaryrd}
% Les diagrammes commutatifs sont partout dans ce fonds : c'est la langue
% même de la géométrie algébrique de Grothendieck, et une transcription qui
% les rendrait en prose ne transcrirait rien.
\usepackage{tikz-cd}
% Français par défaut — c'est la langue des feuillets — mais l'anglais est
% chargé aussi : la traduction et le résumé sont des documents anglais, et la
% typographie française (espaces avant les deux-points) y serait une faute.
% Ces documents basculent par \selectlanguage{english} en tête de corps.
\usepackage[english,french]{babel}
\usepackage{fontspec}
\usepackage{geometry}
\geometry{a4paper,margin=25mm,marginparwidth=28mm}
\usepackage{marginnote}
\usepackage{xcolor}
% ulem, not soul. `soul`'s \st and \ul are built on a character-by-character
% scan that XeLaTeX's font machinery breaks: the document fails with an
% undefined control sequence rather than degrading. ulem does the same two jobs
% and is Unicode-engine safe. `normalem` stops it hijacking \emph, which the
% transcriptions use for Grothendieck's own emphasis.
\usepackage[normalem]{ulem}
\usepackage{hyperref}
\hypersetup{colorlinks=true,linkcolor=black,urlcolor=black,pdfborder={0 0 0}}

% --- Métadonnées du lot -----------------------------------------------------
% Reprises telles quelles par le script de rendu, qui les affiche en tête de
% la vue de lecture. Elles fixent ce que le fichier prétend couvrir : sans
% elles, une transcription flotte sans point d'attache dans un fonds de seize
% mille feuillets.

\newcommand{\folder}[1]{\gdef\grothFolder{#1}}
\newcommand{\batch}[1]{\gdef\grothBatch{#1}}
\newcommand{\leaves}[2]{\gdef\grothFirst{#1}\gdef\grothLast{#2}}
\newcommand{\foldertitle}[1]{\gdef\grothFoldertitle{#1}}
\gdef\grothFolder{?}\gdef\grothBatch{?}\gdef\grothFirst{?}\gdef\grothLast{?}\gdef\grothFoldertitle{}

% --- Le résumé --------------------------------------------------------------
%
% Il ouvre la lecture modernisée, sous le titre « Résumé », et il est composé
% en petit. Ce n'est pas un ornement : c'est le seul passage du document qui
% ne soit pas des mathématiques, et un lecteur doit pouvoir le reconnaître
% avant de l'avoir lu — puis le sauter s'il connaît déjà le sujet, ou s'y
% arrêter s'il ne le connaît pas. Le filet qui le suit marque la reprise du
% corps.

\newenvironment{resume}
  {\small\setlength{\parskip}{0.45em}}
  {\par\medskip\hrule\medskip}

% --- Mots-clés ---------------------------------------------------------------
%
% En anglais, en clôture du résumé de la lecture modernisée : le vocabulaire
% moderne sous lequel chercher ce que les feuillets construisent. Le manifeste
% du site (`npm run manifest`) les extrait pour étiqueter la cote entière —
% cette ligne est la source unique des tags, il n'y a pas de fichier à côté.
\newcommand{\keywords}[1]{\par\smallskip\noindent
  {\footnotesize\textsc{Keywords} --- \textit{#1}}\par}

% --- L'appareil critique ----------------------------------------------------

% Le numéro de feuillet, en marge. C'est l'ancre qui permet de régler un
% désaccord sur une formule en regardant *un* feuillet plutôt qu'une chemise
% de six cents. Le site s'en sert aussi pour tourner les pages du fac-similé
% quand on fait défiler la transcription.
\newcommand{\leaf}[1]{%
  \marginnote{\footnotesize\color{gray!70}\textsf{#1}}%
  \label{leaf:#1}\ignorespaces}

% Illisible : on ne devine pas. Un mot inventé qui se lit comme les autres est
% le pire résultat possible.
\newcommand{\ill}{\textcolor{red!60!black}{[\,\ldots\,]}}

% Lecture douteuse : le mot est donné, et signalé comme incertain.
\newcommand{\uncertain}[1]{\uline{#1}}

% Ajout de l'éditeur — une lettre manquante, un mot sous-entendu.
\newcommand{\add}[1]{\textcolor{blue!50!black}{[#1]}}

% Biffé par Grothendieck lui-même. Ce qu'il a rayé fait partie du document :
% une rature dit souvent où la pensée a changé de direction.
\newcommand{\struck}[1]{\sout{#1}}

% Note du transcripteur, sur l'état matériel du feuillet.
\newcommand{\note}[1]{\par\smallskip{\small\color{gray!60!black}\textit{[#1]}}\par\smallskip}

% Note marginale de Grothendieck, distincte du corps du texte.
\newcommand{\marginal}[1]{\par\smallskip\noindent\hspace{1em}%
  {\small\color{gray!60!black}#1}\par\smallskip}

% --- Titre ------------------------------------------------------------------

\AtBeginDocument{%
  \begin{center}
    {\large\bfseries Fonds Alexandre Grothendieck --- cote n\textsuperscript{o}~\grothFolder}\\[2pt]
    {\itshape\grothFoldertitle}\\[4pt]
    {\small feuillets \grothFirst--\grothLast\ (lot \grothBatch)}\\[2pt]
    {\footnotesize\color{gray!60!black}Transcription établie d'après le fac-similé numérisé
     par l'Université de Montpellier.\\
     Les lectures incertaines sont soulignées, les illisibles notées [\,\ldots\,],
     les ajouts entre crochets.}
  \end{center}
  \vspace{1em}}

\endinput


\folder{43}
\batch{10}
\pages{181}{200}
\dating{[à partir de 1961-vers 1968]}
\watermark{Édition de démonstration}
\foldertitle{Dualité des groupes. Jacobiennes généralisées etc : notes et copies de notes manuscrites (s.d.)}

\begin{document}

\page{181}

$\underline{A}_{\mathfrak{m}}$ \qquad \struck{$\underline{M}$}
\note{en tête du feuillet, deux lettres soulignées d'un trait épais ; la
seconde, un $M$, est biffée. Le petit « $\mathfrak{m}$ » est écrit sous le
$A$, comme les indices des $H^0$ qui suivent ; on le rend par
$\mathfrak{m}$, l'idéal que la fin de la page nomme}

$X$, $Y$ deux modules sur $A$
\[ H^0_{\mathfrak{m}}(X) \xrightarrow{\ \sim\ } D\,H^0_{\mathfrak{m}}(Y) \]

$(X,Y) \to (X',Y')$ consiste en
\[ X \to X', \qquad Y \leftarrow Y' \]
tel que le diagramme

\begin{tikzcd}
H^0(X) \arrow[r] \arrow[d, "\wr"] & H^0(X') \arrow[d, "\wr"] \\
D\,H^0(Y) \arrow[r] & D\,H^0(Y')
\end{tikzcd}

soit commutatif.
\note{dans ce carré l'indice $\mathfrak{m}$ n'est plus écrit ; la flèche du
bas n'a pas de pointe visible et est rendue dans le sens de celle du haut}

Si on veut \uncertain{inclure} les groupes \uncertain{$\mathbb{G}_m$}
\struck{\ill{} \ill{}} \ill{} les \uncertain{algébriques}
(\uncertain{comme aussi} $\mu_p$, $\mathbb{Z}/\ell\mathbb{Z}$) et $\mathbb{Z}$
parmi les \uncertain{groupes} \uncertain{formels} (\uncertain{comme aussi}
$\mathbb{Z}/p\mathbb{Z}$, $\mathbb{Z}/\ell\mathbb{Z}$), alors \ill{} \ill{}
\ill{}, on prend \ill{} \ill{} \ill{} \uncertain{topologique}
\[ \mathbb{Z} \times A \]
et le $H^0$ étant relatif aux idéaux $\mathfrak{m} = (F,V)$ de $A$ et les
idéaux \ill{} \uncertain{quelconques} de $\mathbb{Z}$, la dualité
\note{la phrase se poursuit en haut de la page 182 ; le premier mot après
« Si on veut » et les mots qui précèdent « les algébriques » sont très
rapides et restent en grande partie illisibles}

\page{182}

\uncertain{s'interprète} \ill{} \uncertain{deux} \uncertain{facteurs}.
\note{fin de la phrase commencée page 181 ; le premier mot est très
incertain. Double trait en dessous, qui clôt le passage}

Dans une catégorie de \uncertain{Cartier} \ill{} \ill{} \ill{}, la
\ill{} \ill{} \uncertain{toujours} \ill{} par $F^{-}$, et la \ill{}
\ill{} \ill{} par $V^{-}$. Mais \uncertain{passons} aux ind-objets de $C_1$
et aux pro-objets de $C_2$, les uns ni les autres ne sont
\uncertain{caractérisés} par $F^{-}$ ou par $V^{-}$. Comment définir la
\uncertain{nullité} \uncertain{intrinsèquement} ??
\note{les exposants de $F$ et $V$ sont un petit trait horizontal, rendu
$F^{-}$, $V^{-}$ ; « nullité » peut aussi se lire « nilpotence »}

Quelle est l'interprétation des modules du type fini sur
\[ \mathbb{W}_{\infty\,\sigma}[F,V]/(FV-p) \ ?? \]
\note{le $\sigma$ en indice après $\mathbb{W}_\infty$ est celui des pages
183, 196, 200 : l'anneau non commutatif tordu par le Frobenius}

Je pense que c'est l'anneau des \uncertain{endomorphismes} de
$\mathbb{W}_\infty$ qui \uncertain{laissent} \uncertain{invariant} le
sous-groupe $D(\mathbb{W}_\infty)$ [\struck{\ill{}} \uncertain{groupe}
\ill{} de Witt dont \ill{} les \uncertain{composantes} sont
\uncertain{nilpotentes}, et \uncertain{dont} \ill{} \ill{} de type fini
\ill{} \uncertain{nulles} …], \uncertain{ou} les
\uncertain{endomorphismes} de $D(\mathbb{W}_\infty)$ qui
\uncertain{laissent} \uncertain{invariant} l'accouplement fondamental de
\ill{} \ill{} avec lui-même.

\page{183}
\note{la moitié supérieure de la page (jusqu'au cadre) est barrée de
plusieurs longs traits obliques et d'une courbe ; deux autres traits
obliques traversent le reste de la feuille. Ce qui se lit est donné
ci-dessous}

\struck{$\mathbb{G}_m$ \quad $\mathrm{Ext}$}
\[ 0 \to \mu_n \to \mathbb{G}_m \xrightarrow{\ n\ } \mathbb{G}_m \to 0 \]
\note{le premier terme est surchargé et biffé par l'auteur ; « $\mu_n$ » n'est qu'une
lecture probable, l'indice $n$ étant celui de la suite suivante}
\[ 0 \to \mathbb{Z} \xrightarrow{\ n\ } \mathbb{Z} \to
\mathrm{Ext}^1(\mathbb{G}_m, \mu_n) \to 0 \]

\[ \mathbb{W}_{\sigma}[F,V]/(V^n, FV-p) \]
\struck{$\forall$} \struck{\ill{}} \struck{$k[V]$} \qquad
$k_{\sigma}[F]$ \qquad $k[F]$ \qquad \struck{$F$ \ill{}} $k[F]$

$\widetilde{\mathbb{W}}_\infty$ \uncertain{noethérien},
\uncertain{artinien}, \uncertain{pro-artinien} \ill{}
\uncertain{noethérien}
\note{mots écrits sous $\widetilde{\mathbb{W}}_\infty$ et soulignés, dans la
partie barrée}

\begin{quote}
$\underline{\mathbb{G}_m,\ \mathbb{G}_a}$ \qquad
$\mathbb{Z}/\ell\mathbb{Z}$ --- $\mu_p$ --- $\alpha_p$ ---
$\mathbb{Z}/p\mathbb{Z}$ \qquad
$\mu_p$ --- $\alpha_p$ --- $\mathbb{Z}/p\mathbb{Z}$ ---
$\mathbb{Z}/\ell\mathbb{Z}$ \qquad
$\underline{\mathbb{V}}$ \quad $\underline{\mathbb{Z}}$
\end{quote}
\note{schéma encadré : à gauche $\mathbb{G}_m$, $\mathbb{G}_a$ soulignés ;
deux traits en éventail en partent. Sur le trait supérieur :
$\mathbb{Z}/\ell\mathbb{Z}$, puis $\mu_p$, $\alpha_p$,
$\mathbb{Z}/p\mathbb{Z}$ ; sur le trait inférieur : $\mu_p$, $\alpha_p$,
$\mathbb{Z}/p\mathbb{Z}$, puis $\mathbb{Z}/\ell\mathbb{Z}$ ; les deux
convergent vers $\mathbb{V}$ et $\mathbb{Z}$ (lettres doublées) soulignés,
à droite. Le schéma met face à face, dans l'ordre inverse, les groupes que
la page 192 énumère}

\[ \mathbb{W}_{\sigma}\{F,V\}/(FV-p) \longrightarrow
\mathbb{W}_{\sigma}[[F,V]]/(FV-p) \]
(\uncertain{en} dimension 2)

\emph{Anneau de Cartier} \qquad \emph{Anneau de Gabriel-Dieudonné}
\note{« Anneau de Cartier » est souligné deux fois ; « Anneau de
Gabriel-Dieudonné » est souligné d'une flèche et d'un trait. Entre les
deux, une insertion entourée : « (\uncertain{pas} \uncertain{nécessairement}
\struck{\ill{}}) »}

\marginal{Unifier les pro-groupes algébriques
\uncertain{unipotents} \struck{\ill{}} et les groupes \struck{\ill{}}
(ind-finis)}

\uncertain{complétion} des précédents pour l'idéal \add{maximal}
$\mathfrak{T}(F,V)$
\note{« maximal » est ajouté au-dessus de la ligne avec un renvoi. La
lettre devant $(F,V)$ est une grande lettre ronde, rendue
$\mathfrak{T}$ ; lecture incertaine}

Classification des
\note{la page s'arrête sur ces mots}

\page{184}

\[ \mathbb{W}_\infty \to \mathbb{W}_\infty
 = \varinjlim \mathrm{Hom}(\mathbb{W}_\infty, \mathbb{W}_n) \]
\note{le premier membre est écrit sans nom de flèche ; on lit
$\mathrm{Hom}(\mathbb{W}_\infty, \mathbb{W}_\infty)$ comme limite inductive
des $\mathrm{Hom}(\mathbb{W}_\infty, \mathbb{W}_n)$, l'indice du second
$\mathbb{W}$ étant très abrégé}
\[ D(\mathbb{W}_\infty) \subset \mathbb{W}_{\infty 1} \]
\note{le $D$ est écrit au-dessus d'un « Hom » biffé ; en dessous, un
$D$ surchargé et biffé. L'indice « 1 » de $\mathbb{W}_{\infty}$ est net}

$F$
\[ D(\mathbb{W}_\infty), \qquad D(\mathbb{W}_\infty) \longrightarrow
\mathbb{W}_\infty, \qquad \underline{D(\mathbb{W}_\infty)} \times
\underline{D(\mathbb{W}_\infty)} \]
\[ \mathrm{Hom}(\mathbb{W}_\infty, \mathbb{G}_m), \qquad \mathbb{W}_\infty \times \]
\note{calculs épars dans la moitié droite du haut de la page ; le
« $\mathbb{W}_\infty \times$ » reste en suspens}

\begin{quote}
groupes algébriques \emph{affines} $=$ algèbres de type fini sur $k$,
\uncertain{avec} $\Delta$ \uncertain{commutative}
\uncertain{fonctorielle} / \ill{}

\struck{groupes} \struck{\ill{}} pro(algébriques affines) $=$
\uncertain{Algèbres} \ill{}, \ill{} type fini

\emph{Question} : est-ce un \ill{} pro(groupe alg.\ affine), i.e.
l'algèbre est-elle limite inductive d'algèbres de type fini qui sont
stables par $\Delta$ ??
\end{quote}
\note{les trois énoncés sont réunis par une grande accolade à gauche ;
les deux lignes de la \emph{Question} par une seconde accolade. Après
« avec $\Delta$ », deux mots superposés entre accolades, le premier lu
« fonctorielle », le second illisible}

\page{185}

\struck{hyperalgèbres} \struck{groupes} \add{objet} (ind(algébrique
fini)) $\simeq$ coalgèbres \uncertain{commutatives} \ill{}
\uncertain{unitaires}, \ill{} \uncertain{engendrées} \uncertain{par}
\uncertain{ses} \uncertain{sous-coalgèbres} de \uncertain{rang} fini sur
$k$
\note{sous « $\simeq$ », une colonne de mots serrés :
« \ill{} \ill{} \ill{} en corps, et \ill{} pas utiliser la dualité » ;
lecture très incertaine. « objet » est écrit sous les deux mots biffés}

[\struck{Toute telle coalgèbre est en effet} \uncertain{engendrée} par des
coalgèbres (cette dernière condition \ill{} \ill{}
\uncertain{supplémentaires} \ill{} \ill{}, i.e.\ une multiplication
(\emph{continue} sur $k^I$ n'est pas \ill{}) \ill{} \ill{} \ill{} \ill{}
\uncertain{projective} d'algèbres de \uncertain{rang} fini sur $k$)
\note{le crochet ouvrant et le début de la première ligne sont biffés
d'un trait horizontal}

\marginal{$U \cdot U$ \quad $\Delta(x) + U \otimes U)$}
\note{en marge gauche, un petit schéma : $U$, $U$ en haut, $U^{*}$,
$U^{*}$ en bas, reliés par deux flèches croisées}

\struck{\ill{}}

groupe (ind(algébrique fini)) $\simeq$ coalgèbres de type \ill{},
munies d'une loi d'algèbre \ill{} commutative, unitaire,
\uncertain{compatible} \ill{} les \uncertain{suivantes}
\note{« groupe (ind(algébrique fini)) » est souligné d'un long trait ;
sous lui, entouré : « \uncertain{groupe} \uncertain{formel} »}

\emph{Exemple} : l'algèbre affine d'un groupe algébrique linéaire
\uncertain{ou} \ill{} d'un pro(groupe algébrique linéaire) \ill{} :
d'après Cartier, \add{(\uncertain{en} car.\ 0)} une telle algèbre serait
limite inductive d'algèbres de type fini \ill{} \ill{} \ill{}
\uncertain{de} type fini du \ill{} \ill{} \uncertain{associatif}
[\uncertain{du} \uncertain{moins} \ill{} \ill{} \ill{} alg.\ \ill{},
\uncertain{si} la partie discrète \ill{} \uncertain{réduite} à la
partie infinitésimale]
\note{la phrase de Cartier est écrite dans la colonne de droite, séparée
de « Exemple » par un grand trait courbe ; « (en car.\ 0) » est inséré
au-dessus avec un renvoi, et une ligne d'insertion entourée (« \ill{}
type, \ill{} \ill{} ») la précède}

Donc il en résulte une dualité entre
\begin{quote}
pro(groupes algébriques linéaires) $\sim$ gr(ind.\ alg.\ finis).
\end{quote}
Le plus beau serait manifestement qu'on ait le résultat que cet
\uncertain{anneau} \ill{} \ill{} \uncertain{hyperalgèbres}
\uncertain{analogues}.
\note{la page s'arrête ici ; le dernier mot est souligné}

\page{186}

La catégorie \uncertain{des} \struck{\ill{}} groupes \uncertain{formels}
\ill{} \uncertain{sur} un groupe pro-algébrique linéaire / \uncertain{définissons}
dans une \struck{catégorie} hyperalgèbre topologique \uncertain{pure}
\ill{} $k$, \uncertain{dont} l'espace \uncertain{sous-jacent} est
linéairement compact. Il n'est pas exclu que les hyperalgèbres
topologiques \uncertain{ayant} cette propriété \struck{\ill{}} puissent
être interprétées ainsi. \struck{\ill{}}
\note{« / définissons » est relié par un trait oblique à
« hyperalgèbre topologique », ligne suivante ; le dernier mot du
paragraphe est raturé en zigzag}

\marginal{(\uncertain{définis} comme \uncertain{catégorie}
\uncertain{pro} (ind(\uncertain{$k$-affines}))) \ill{} \ill{}}
\note{en marge gauche, écrit en oblique, un petit bloc serré dans un
cadre partiellement raturé ; seuls les premiers mots se lisent}

Le mécanisme de dualité de Cartier est alors immédiat, ainsi que son
interprétation en termes d'accouplement dans $\mathbb{G}_m$. La
définition de $F$ et $V$ en \uncertain{suit} immédiatement, en car.~$p$.

\[
\begin{array}{ccc}
\mathbb{G}_m & \mathbb{Z}/\ell & \mu_p \\[4pt]
F(x) = x^p & F(x) = x^{?} & F = 0 \\
V(x) = x & V(x) = x^p & V = \mathrm{id}
\end{array}
\]
\note{tableau au bas de la page ; une longue courbe relie
$\mathbb{G}_m$ à $\mathbb{Z}/\ell$, une autre traverse la colonne du
milieu. L'exposant de $F(x) = x^{\cdot}$ sous $\mathbb{Z}/\ell$ est
noyé sous une surcharge (rendu $?$) ; sous la colonne,
« ($F$ et $V$ \uncertain{bijectifs}) ». Au-dessus de $F(x) = x^p$ dans
la première colonne, un « $F = x^p$ » biffé. Deux flèches mènent de $\mathbb{G}_m$ et de
$\mathbb{Z}/\ell$ vers une note en oblique, en bas à gauche :
« \ill{} la \uncertain{réduction} \ill{} \ill{} : $\mathbb{F}_p$
\ill{} \uncertain{identité} \ill{}
$\mathbb{G}_a = \mathbb{G}_a^{(p)}$ »}

\page{187}

$A$ algèbre affine d'un groupe affine $G$ (\uncertain{pas}
\uncertain{nécessairement} de type fini) sur l'anneau $k$
\begin{quote}
$G \to \mathbb{G}_m$ homom.\ de $k$-groupes $\approx$ $G \to \mathbb{G}_a$
morphismes de $k$-schémas,
\end{quote}
tels que (i) Addition transformée en multiplication, (ii) $0$ s'envoie
sur $1$
\note{« homom.\ de $k$-groupes » est écrit sous la première flèche,
« morphismes de $k$-schémas, tels que (i) … (ii) … » en colonne à
droite de la seconde. Le $\mathbb{G}_a$ est tracé en capitale grasse}

i.e.\ $f \in A$, telle que
\begin{quote}
a) $\Delta f = f \otimes f$

b) $\varepsilon(f) = 1$
\end{quote}

Supposons que $A$ soit \emph{libre}, alors $A$ est l'ensemble des formes
$k$-linéaires continues sur $\check{A} = \mathrm{Hom}(A,k)$, muni de la
convergence simple. Alors, a) signifie que $f$ est multiplicative
\struck{\ill{}} [quand $\check{A}$ est muni de la structure de $k$-algèbre
déduite de $\Delta$] et b) que $f$ est unitaire ; donc a) et b), que $f$
est un homomorphisme d'algèbres topologiques. \struck{D'où, si} Notons que
rien ne garantit que $\check{A}$ soit topologisé linéairement
\uncertain{par} \ill{} \uncertain{idéaux}, i.e.\ qu'il y ait un système
fondamental de voisinages qui soient des \emph{idéaux} ; or, ce qui
revient au même \ill{} si $k$ est un corps [et ce qui \uncertain{semble}
\ill{} \ill{} \ill{}]. \uncertain{rien} \uncertain{ne}
\note{la phrase continue en haut de la page 188, « garantit que $A$ … »}

\page{188}

garantit que $A$, en tant que coalgèbre, soit \uncertain{réunion} de
« sous-coalgèbres » de type fini \struck{l'est, de façon} en tant que
modules sur $k$\struck{. S'il en est ainsi,} [et à fortiori
\uncertain{plusieurs} \uncertain{difficultés} ne \ill{} pas
\uncertain{rencontrées} \ill{} $k$]. \emph{S'il} en est
\uncertain{cependant} ainsi, alors $\check{A}$ peut être regardé comme
l'algèbre affine d'un ind-(schéma affine) sur $k$, [\uncertain{définissant}
un schéma formel] \struck{\ill{} \ill{} groupe} \struck{\ill{} \ill{}
\ill{}} et \struck{\ill{}} \ill{} \ill{} un \struck{groupe} (ind-affine)
sur $k$, $D(G)$ et appelé \emph{groupe dual} de $G$. \struck{La
\ill{} \ill{}} La \uncertain{construction} \uncertain{précédente}
\struck{\ill{}} \uncertain{est} \ill{} \uncertain{une}
\uncertain{construction} alors par \ill{} \struck{\ill{}} \uncertain{un}
groupe ind.\ affine (\uncertain{au} sens Cartier) \ill{} \ill{} de
$k$-algèbres (\uncertain{représentable} par un \ill{}) \uncertain{non}
\uncertain{précisément} \ill{} $\underline{\mathrm{Hom}}_{k\text{-gr.}}(G,
\mathbb{G}_m)$.
\note{moitié haute de la page très rapide, avec plusieurs ratures et
insertions interlinéaires ; la marge gauche porte, en oblique : « Si,
\uncertain{dans} \ill{} \ill{} \ill{} \ill{} \uncertain{bien}… », qui
renvoie par des points au crochet « [et à fortiori » ; lecture très
incertaine}

Si $G_k$ est \struck{de type fini} un \ill{} de $G$ de type fini sur $k$,
alors si $k$ est de car.\ $p > 0$, d'après Cartier $D(G)$ est bien un
groupe (ind-fini),
\note{l'indice de $G_k$ est surchargé ; le mot après « un » reste illisible}
\[
\left.\begin{array}{l}
D(\mathbb{G}_m) = \mathbb{Z} \\
D(\mathbb{Z}/\ell\mathbb{Z}) \simeq \mathbb{Z}/\ell\mathbb{Z} \\
D(\mathbb{Z}/p^n\mathbb{Z}) \simeq \mu_{p^n} \\
D(\mu_{p^n}) \simeq \mathbb{Z}/p^n\mathbb{Z} \\
D(\mathbb{G}_a) = \ \cdots
\end{array}\right.
\]
\note{les cinq formules sont réunies par une accolade. Au-dessus, un mot
raturé. En face de $D(\mathbb{Z}/\ell\mathbb{Z})$ : « \struck{\ill{}}
\uncertain{si} $\ell$ premier \ill{} car., et si \ill{} \ill{}
\uncertain{par} \ill{} \ill{} \ill{} de $k$ » ; en face de
$D(\mathbb{Z}/p^n\mathbb{Z})$, un mot biffé. Pour $D(\mathbb{G}_a)$ :
« $=$ \struck{algèbre} \ill{} \ill{} contenant de la $p$-algèbre de Lie
libre engendrée par un élément (\uncertain{même} un corps de car.\
$p > 0$) », avec en insertion au-dessus « \ill{} \uncertain{définition}
\ill{} \ill{} » ; la phrase s'arrête au bas de la page}

\page{189}

Soit $\underline{C}$ une catégorie, \uncertain{d'où}
\[ \mathrm{Pro}(\underline{C}) \subset
\underline{\mathrm{Hom}}(\underline{C}, (\mathrm{Ens})) \]

Un objet $\mathfrak{X}$ de $\mathrm{Pro}(\underline{C})$ est
\uncertain{considéré} \uncertain{toujours} \uncertain{comme} (\uncertain{par}
\uncertain{abus}) équivalent à \ill{} \uncertain{foncteur}
\[ Y \mapsto \mathrm{Hom}_{\mathrm{Pro}(\underline{C})}(Y, \mathfrak{X}) \]
de $\mathrm{Pro}(\underline{C})^{\circ}$ \uncertain{dans} $(\mathrm{Ens})$,
\struck{\ill{}} \uncertain{mais},
\[ \mathfrak{Y} = (Y_j)_{j\in J}, \qquad
\mathrm{Hom}_{\mathrm{Pro}(\underline{C})}(\mathfrak{Y}, \mathfrak{X})
= \varprojlim_{i} \mathrm{Hom}(\mathfrak{Y}, \mathfrak{X}_i) = \]
\note{cette ligne est en partie biffée d'un trait (« $\mathfrak{Y} =
(Y_j)_{j\in J}$, les $Y_i \in \underline{C}$, \ill{} \ill{} ») et la
formule est reliée par un trait à la ligne suivante, qu'elle remplace ; le
second membre reste sans suite}

identifions $\underline{C}$ à une sous-catégorie de
$\mathrm{Pro}(\underline{C})$, on notera que $\mathfrak{X}$ \uncertain{est}
\emph{\uncertain{plus} \uncertain{commode}} : \uncertain{il} \ill{}
puis \uncertain{prendre} \ill{} \uncertain{restreint} le foncteur
$Y \mapsto \mathrm{Hom}_{\mathrm{Pro}(\underline{C})}(Y, \mathfrak{X})$ à
$\underline{C}^{\circ}$ dans $(\mathrm{Ens})$. [\emph{Ex} :
$\underline{C} = (\mathrm{Ens})$, alors
$(X_i)$ et $X = \varprojlim_{(\mathrm{Ens})} X_i$ définissent le
\uncertain{même} foncteur
$\underline{C}^{\circ} \to (\mathrm{Ens})$, \uncertain{mais} \ill{}
\uncertain{qui} \uncertain{ne} \uncertain{sont} \uncertain{pas}
\uncertain{isomorphes} en $\mathrm{Pro}(\underline{C})$]. Le foncteur
\uncertain{naturel}
\[ \mathrm{Pro}(\underline{C}) \longrightarrow
\underline{\mathrm{Hom}}(\underline{C}^{\circ}, (\mathrm{Ens})) \]
\uncertain{n'est} donc pas pleinement
fidèle en général. Il \uncertain{importe} de \uncertain{bien}
distinguer \uncertain{en} particulier, \uncertain{si}
$X = \varprojlim_{\underline{C}} \mathfrak{X}$ existe, entre
$\mathfrak{X}$ et $X$ [\uncertain{on} \uncertain{peut} \uncertain{dire}
\uncertain{semble-t-il} \uncertain{que} \ill{} \ill{} \ill{}
\uncertain{naturel} \uncertain{de} $\mathrm{Pro}\,\underline{C}$
\[ X = \varprojlim_{\underline{C}} \mathfrak{X} \longrightarrow \mathfrak{X} \]
\note{« donc » est écrit au-dessus de la ligne, avec une accolade, sur un mot
biffé et souligné. La page s'arrête sur cette flèche, le crochet
n'étant pas refermé}

\page{190}

universel par les morphismes \add{$\mathrm{Pro}(\underline{C})$} de
l'objet $X$ dans $\mathfrak{X}$ ; si $\varprojlim_{\underline{C}} \mathfrak{X}$
existe pour tout $\mathfrak{X}$, cela signifie précisément que le foncteur
\struck{\ill{}}
\[ \underline{C} \longrightarrow \mathrm{Pro}(\underline{C}) \]
\uncertain{canonique} \uncertain{pleinement} fidèle \add{\uncertain{admet}}
\struck{\ill{}} \uncertain{un} \uncertain{adjoint}, « \uncertain{projection} »
de $\mathrm{Pro}(\underline{C})$ \struck{\ill{}} sur $\underline{C}$.
\note{« $\mathrm{Pro}(\underline{C})$ » est ajouté au-dessus de « morphismes »
avec une flèche ; un mot est inséré au-dessus de « \ill{} de $K$ »,
à la troisième ligne. Lecture incertaine de toute cette phrase}

\uncertain{Comme} \ill{} \ill{} \uncertain{conditions} pour cette
situation \ill{} la \uncertain{suivante} : \struck{\ill{}} il \ill{}
\uncertain{existe} pour \uncertain{toute} \uncertain{limite} projective
de $\underline{C}$, \uncertain{mais} \ill{} \ill{} \ill{}
\emph{sous-catégorie pleine} \ill{} de $\underline{C}'$ \uncertain{telle}
que \struck{\ill{}} \uncertain{tout} système projectif \uncertain{dans}
$\underline{C}$ ait une limite projective dans \emph{$\underline{C}'$} ;
\uncertain{donc} \ill{} \ill{} \uncertain{un} foncteur naturel
\[ \mathrm{Pro}(\underline{C}) \longrightarrow \underline{C}' \]
qui \emph{\uncertain{n'est}} \struck{\ill{}} (\ill{} \ill{} \ill{}
\uncertain{triviaux}) \uncertain{pas} pleinement fidèle. [\ill{} \ill{}
\ill{} p.\ ex.\ \uncertain{naturel} de \uncertain{quels} \ill{} de
\ill{} \ill{}, pris \uncertain{sous} \ill{} \uncertain{limite} de
$\underline{C}$, \uncertain{sont} \uncertain{artiniens} ??]

\emph{Exemple} : les pro-\uncertain{schémas} affines de type fini
\uncertain{sur} $S$ \struck{\ill{}} (\struck{\ill{}} \ill{} \ill{}
\ill{} \ill{}) \uncertain{préschéma} loc.\ \uncertain{noethérien})
s'identifient aux schémas \struck{affines} affines sur $S$.

\emph{Question} \struck{\ill{} $S$ \ill{} \ill{}} Un schéma en groupes
affine sur $S$ est-il un pro(groupe affine de type fini), i.e.\ si
$S = \mathrm{Spec}(k)$
\note{la page s'arrête sur cette hypothèse}

\marginal{N.B.\ $\mathrm{Pro}(\underline{C}\text{-}\mathrm{Gr}.)$ \ill{} $\to (\mathrm{Pro}\,\underline{C})\text{-}\mathrm{Gr}.$ \ill{} \uncertain{pleinement} fidèle}
\note{en marge gauche, en oblique, avec une flèche ; la flèche
elle-même et le mot « pleinement » sont de lecture incertaine}

\page{191}

(\uncertain{k} anneau \uncertain{quelconque}) est-il vrai que si $A$
\uncertain{est} une hyperalgèbre (\ill{} \ill{} \ill{} \ill{} ?) sur
$k$, alors $A$ \struck{\ill{}} est limite inductive de
\struck{\ill{}} \add{\uncertain{sous-hyperalgèbres}} \uncertain{sous-algèbres}
de type fini \add{\uncertain{en} \uncertain{tant} \uncertain{que}
\uncertain{module}} \struck{\ill{} \ill{}}. Il suffit \uncertain{évidemment}
que $A$ soit limite inductive de \uncertain{sous-$k$-modules} de type fini
$V$ tels que $\Delta V \subset \mathrm{Im}\, V \otimes V$, et tels que
[$A_V$ \uncertain{étant} la sous-$k$-algèbre \uncertain{engendrée}
\uncertain{par} \ill{}] \struck{\ill{}}
$A_V \otimes A_V \to A \otimes A$ soit \emph{injectif}. \struck{\ill{}}
\note{les deux insertions sont écrites au-dessus de la ligne, au-dessus
de mots biffés ; « en tant que module » est de lecture très incertaine}

\emph{\uncertain{Cela} \uncertain{est} \uncertain{facile} \uncertain{si}
\ill{} \ill{} \ill{} \ill{}} $k$ est un \emph{corps} [\uncertain{dans}
le \struck{\ill{}} \add{\uncertain{deuxième}} \uncertain{cas} \ill{}
\ill{} ; \uncertain{pour} le \uncertain{premier}, \ill{} \ill{} \ill{}
\struck{\ill{}} \add{\ill{}}] ; \uncertain{et} \uncertain{toujours}
\ill{} \uncertain{possible} \ill{} \uncertain{si} $k$ est
\emph{\uncertain{noethérien}}, \uncertain{condition} \ill{} \ill{}
\uncertain{plus} \uncertain{générales}, \uncertain{telles} que :
$A$ \uncertain{libre}].
\note{passage très rapide ; les mots soulignés en tête sont peu lisibles.
Au bord droit, un petit « $k_y$ » sans rattachement}

\emph{Donc} \uncertain{toute} \uncertain{hyperalgèbre} \uncertain{sur}
$k$, \struck{\ill{} \ill{} affine}

$\underline{\text{pro(groupe alg.\ affine)}} =
\underline{\text{groupe alg.\ affine}}$
\note{deux lignes soulignées en tête du paragraphe, suivies d'un signe
raturé}

Le « dual » d'un \struck{\ill{}} groupe \uncertain{affine}
\add{\uncertain{commutatif} \uncertain{sur} $k$} \uncertain{est}
\uncertain{ensuite} \ill{}, condition (1) est un groupe (ind-fini)
$\hat{D}(G)$, \uncertain{à} \uncertain{Cartier} \uncertain{sur} $S$, dont
l'hyperalgèbre est l'algèbre affine de $G$. \uncertain{La}
\uncertain{condition} \ill{} type \uncertain{fini} \uncertain{sur}
\uncertain{un} \uncertain{corps} alg.\ clos, \ill{} \ill{}
\uncertain{comme}
\note{« commutatif sur $k$ » est inséré au-dessus de la ligne. Le
circonflexe de $\hat{D}(G)$ est posé au-dessus de $G$ seul ; la graphie
de $D$ est surchargée. En marge gauche, en oblique : « \uncertain{Voir}
\uncertain{aussi} \ill{} \ill{} Cartier ». La page s'arrête en milieu de
phrase}

\page{192}

\uncertain{types} \struck{\uncertain{simples}} \add{« \uncertain{élémentaires} »},
de \uncertain{groupes} \uncertain{dual} \uncertain{l'un} de l'autre\,:
\note{« élémentaires » est écrit au-dessus d'un mot biffé}
\[
\begin{array}{ll}
\mathbb{G}_m,\ \mathbb{Z}/\ell\mathbb{Z},\ \mu_p\ \ \mathbb{G}_a,\ \alpha_p,\ \mathbb{Z}/p\mathbb{Z}
 & \qquad \mu_p,\ \alpha_p,\ \mathbb{Z}/p\mathbb{Z} \\[4pt]
\mu_p,\ \alpha_p,\ \mathbb{V},\ \mathbb{Z}/p\mathbb{Z},\ \mathbb{Z}/\ell\mathbb{Z},\ \mathbb{Z} &
\end{array}
\]
\note{les deux listes sont l'une sous l'autre, la seconde renversant la
première ; à droite de la première, « $\mu_p, \alpha_p, \mathbb{Z}/p\mathbb{Z}$ »
avec en dessous « \uncertain{spéciaux} : \ill{} pro » }

[\uncertain{Ceux-ci} ; l'un \ill{} \ill{} les \uncertain{« éléments »}
\ill{} \ill{}, \ill{} de \uncertain{composition} \uncertain{complète}
\ill{} \ill{}, \uncertain{i.e.} \ill{} \ill{} \ill{} \ill{} \ill{}
\ill{} \uncertain{premiers}, \uncertain{dans} \ill{} \ill{} \ill{}
alg.\ clos \ill{} $=$ \uncertain{v-biais} ]
\[ \mathbb{V} = D(\mathbb{G}_a) \]
est le \uncertain{noyau} \ill{} \ill{} ; \ill{} \ill{} \ill{}, les
\uncertain{seuls} \uncertain{ensemble} \ill{} \uncertain{Cela} \ill{}
\uncertain{considère} \ill{} \uncertain{les} \ill{} faisant des
\emph{extensions}, \uncertain{Produits} \struck{\ill{} \ill{}}, \ill{}
\uncertain{catégorie} (\uncertain{pro}\uncertain{-noethérienne}
\uncertain{artinienne}) des groupes (ind-affines)
\struck{[\uncertain{Classification}} [\uncertain{ind}-\ill{} \ill{}
\ill{} \ill{} des \uncertain{éléments} \uncertain{dessus}] \ill{} \ill{}
\uncertain{une} algèbre affine, \add{\uncertain{définie}} qui \ill{}
\ill{} \uncertain{directe} \ill{} \ill{} \ill{} \ill{} \struck{\ill{}
\uncertain{topologique} \ill{}} \ill{} \ill{} \ill{} \uncertain{dimension}
(\uncertain{Borel} \ill{}) \ill{} (\ill{} \ill{} \ill{}) \ill{}
\uncertain{d'une} \ill{} \ill{} \ill{} \uncertain{solution} \ill{}
\uncertain{pour} (\uncertain{dans} le cas \ill{} \uncertain{géométrique})
\note{moitié médiane de la page très rapide, surchargée d'insertions ;
seuls les termes mathématiques se lisent avec quelque sûreté}

\struck{Définition de $F$, $V$,}

\begin{quote}
\emph{Multiplicatif} [\uncertain{quasi-compl.} \ill{} \uncertain{communs}]
\quad \emph{\uncertain{unipotent}} \quad [\uncertain{sans}
\uncertain{l-torsion} \uncertain{sur} \ill{}] \uncertain{groupes}
\uncertain{discrets}

\struck{Tore} \struck{\ill{}} groupe profini : \struck{\ill{}} $\mu_p$
\qquad \ill{} \ill{} \uncertain{unipotent} \ill{} \ill{} \uncertain{alg.}
\ill{} \ill{} \uncertain{profini} \uncertain{(\ill{})} \qquad
\uncertain{réduit} \ill{} \ill{} $\pi$ \ill{} \uncertain{continu}
\end{quote}
\note{bas de la page : une petite table en trois colonnes, reliées par
des traits horizontaux à bornes, sous les en-têtes « Multiplicatif »,
« unipotent » (au-dessus d'un mot biffé) et « groupes discrets » ; les
cases, en écriture minuscule et raturées, portent notamment « Tore »,
« groupe profini », $\mu_p$ (avec une grande croix devant), et plusieurs
fois la lettre $\pi$. Lecture très partielle}

\marginal{\uncertain{Sur} \ill{} \ill{} \uncertain{universelle}
\ill{} \ill{} \ill{} \ill{} \uncertain{extension} \ill{} \ill{}
\uncertain{Dieudonné} \ill{} \struck{\ill{}} \uncertain{hyperalgèbre}
\ill{} \ill{} \ill{} \ill{} \uncertain{quelconque} \ill{}
\ill{} \uncertain{Cartier} \ill{} \ill{} \uncertain{V-groupe}
\uncertain{hyperfini} \ill{}}
\note{deux longues notes en oblique occupent toute la marge gauche,
du haut au bas de la page ; elles restent presque entièrement
illisibles}

\section*{Catégories abéliennes}
\note{titre souligné, en tête de la page 194 ; la page change d'encre et
de papier, et l'écriture devient posée jusqu'à la page 200}

\page{194}

Soit $\underline{C}$ une catégorie abélienne où tous les objets sont de
longueur finie. Alors $\mathrm{Ind}\,\underline{C}$ [catégorie des
foncteurs contravariants exacts à gauche de $\underline{C}$ dans
\struck{\ill{}} $(\mathrm{Ab})$] est une \emph{catégorie abélienne} à
\emph{limites inductives exactes}, à générateurs de longueur finie
[\uncertain{nous} \uncertain{dirons} : \emph{loc.\ de longueur finie}]
dont $\underline{C}$ est la sous-catégorie des objets de longueur finie.
Il y a ainsi essentiellement correspondance biunivoque entre catégories
abéliennes à objets de longueur finie, et catégories abéliennes à limites
inductives exactes à générateurs de longueur finie.

Soit $\underline{C}$ une catégorie abélienne \add{(à limites inductives
exactes)}, $X \in \underline{C}$, le \emph{socle} de $X$ est
\struck{la} somme des sous-objets simples de $X$, soit $F_0(X)$. On
désigne par $F_1(X)$ l'image inverse du socle de $X/X_0$ etc., et on
trouve ainsi une filtration croissante. \struck{\ill{}}
\note{« (à limites inductives exactes) » est écrit au-dessus de la ligne
avec une accolade de renvoi. Dans « l'image inverse du socle de
$X/X_0$ », $X_0$ désigne $F_0(X)$}

\page{195}

Cette \emph{filtration est fonctorielle}. $\underline{C}$ loc.\ de
longueur finie $\Longleftrightarrow$ tout $X$ est la sup des $F_i(X)$.
Soit en général $\Phi_0(X) = \mathrm{Sup}\, F_i(X)$, c'est le plus grand
sous-objet de $X$ qui est « de dimension de Krull $0$ ». C'est le noyau
de $X \to X_{\underline{C}_0}$ (localisé par $\underline{C}_0$) où
$X_{\underline{C}_0}$ est le localisé de $X$ pour la sous-catégorie
$\underline{C}_0$ de $\underline{C}$ des objets \uncertain{introduite}
\uncertain{dans} \uncertain{de} \uncertain{dimension} \uncertain{de}
Krull $0$. $\Phi_1(X)$ \uncertain{sera} l'image inverse de
$\Phi_0(X_{\underline{C}_0})$ \struck{\ill{}} [\uncertain{ou}
\uncertain{encore} comme $\mathrm{Ker}(X \to X_{\underline{C}_1})$]
etc.\ \uncertain{On} \uncertain{trouve} \uncertain{ainsi} une
\emph{filtration fonctorielle} $(\Phi_i(X))_i$. Dire que
$\underline{C}$ est de dimension de Krull $\leq n$ signifie que
$\Phi_n(X) = X$ pour tout $X$.
\note{la lettre notée $\Phi$ est un $\Phi$ de forme ronde, un peu
gothique ; on la rend par $\Phi$. Le passage « où $X_{\underline{C}_0}$
est le localisé de $X$ pour la sous-catégorie $\underline{C}_0$ de
$\underline{C}$ des objets » est un ajout interlinéaire, relié par une
accolade à « \uncertain{introduite} dans » ; en marge gauche, à la même
hauteur : « \uncertain{de} \uncertain{dimension} \uncertain{de} Krull $0$ ».
Dans $\Phi_0(X_{\underline{C}_0})$ l'indice $\underline{C}_0$ est écrit sous
un mot noirci. « [ou encore comme $\mathrm{Ker}(X \to
X_{\underline{C}_1})$] » est inséré au-dessus de la ligne, avec un
renvoi}

\page{196}

\struck{\ill{}}

ind(pro-(groupes alg.\ finis) \add{\uncertain{unipotents}}
\struck{\uncertain{artiniens}}) $=$ pro(Modules loc.\ de long.\ finie,
\uncertain{ind} \uncertain{artiniens}) dual des Modules sur $A$.
\note{« unipotents » est écrit au-dessus d'un mot biffé, lu
« artiniens »}

\begin{tikzcd}[column sep=small, row sep=large, nodes={font=\scriptsize}]
\text{\emph{groupes formels raisonnables}} \arrow[r, leftrightarrow, "\text{contrav.}"] \arrow[d, leftrightarrow, "\text{contrav.}"] & \text{Modules \ldots\ sur } A \\
\text{\emph{pro-(groupes alg.\ finis) raisonnables}} \arrow[ur, leftrightarrow, "\text{cov.}"'] &
\end{tikzcd}

\note{schéma redessiné : « groupes formels raisonnables » (souligné) est
relié par une double flèche marquée « contrav. » aux « Modules
\struck{de type} finis sur $A$ tels que $M/FM$ soit de longueur finie »,
et par une double flèche verticale, marquée d'un mot lu « contrav. »,
aux « pro-(groupes alg.\ finis) raisonnables » ; de ceux-ci une
troisième double flèche, marquée « cov. », remonte vers la colonne de
droite. Les étiquettes « contrav. » et « cov. » sont des lectures
incertaines ; un signe raturé se trouve sous « de longueur finie »}

\struck{$F\ V$}
\[ \sum_{-\infty}^{+\infty} w_i F^i + \sum_{j}^{\infty} w_j V^j \]
\note{devant la somme, un mot biffé ; au-dessus, un signe raturé et
« $F\ V$ » biffés}

\uncertain{$w_i = 0$} si $i \to \infty$ \qquad
$W_{\sigma}[F][F^{-1}]$
\note{un grand trait oblique sépare ces deux inscriptions ; à droite,
trois lignes biffées de traits obliques (« \ill{} \ill{} \ill{} \ill{}
\ill{} \uncertain{graduation} \ill{} ») sont barrées}

\[
\begin{aligned}
F^{-1} w_i &= \sigma^{-1}(w_i)\, F^{-1} \\
F^{-1} w_j &= \sigma^{-1}(w_j)\, F^{-1} \\
w_i F &= F \sigma^{-1}(w_i)
\end{aligned}
\qquad\qquad
\begin{aligned}
FV &= p \\
V &= pF^{-1}
\end{aligned}
\]
\note{sous « $V = pF^{-1}$ », une formule encadrée « $F^{-1}V = $ » et
une autre, « $F^{-1}V$ » (avec un trait sur $F$), sont biffées de
traits obliques. À droite, un signe raturé et « $FV = p$ » répété.
Dans la troisième ligne, la parenthèse ouvrante manque sur la page et
le $F$ de droite est écrit en minuscule}

\page{198}
\note{la page est écrite sur le recto d'une feuille à en-tête imprimé
(en haut à gauche), dont la page 197 est le verso vierge. Les $W$ de
cette page sont écrits d'un seul trait, non doublés, sauf dans
$\mathbf{W}^{\xi}$ ; on les garde tels, sans les confondre avec les
$\mathbb{W}$ des pages 182--184 et 196}

\[ W_{n,k} = \mathrm{Ker}\bigl(W_n \xrightarrow{\ F^k\ } W_n\bigr) \]

\begin{tikzcd}
W_{n,k} \arrow[r, "i"] \arrow[d, "V"'] & W_{n,k'} \arrow[d] \\
W_{n+1,k} \arrow[r, "i"] & W
\end{tikzcd}

$k' \geq k$
\note{le coin inférieur droit porte $W$ sans indices}

\[ \overline{W}_n = \bigcup_k W_{n,k} \qquad
W_{1,k} \overset{V}{\hookrightarrow} W_{2,k} \hookrightarrow W_{3,k}
\hookrightarrow \cdots \]
\[ \overline{W}_1 \to \overline{W}_2 \to \cdots \overline{W}_n \to \cdots
\qquad \mathcal{U} = \varinjlim_i \overline{W}_i \]

\marginal{\uncertain{est-ce} une \uncertain{solution} de $A$ ??}
\note{inscription au crayon, en oblique, entourée d'un ovale, à gauche}

\[ F : \mathcal{U} \to \mathcal{U}, \qquad V \]

\begin{tikzcd}
W_4 \arrow[r, no head, "V_{4,5}"] & W_5 \arrow[r, no head] & W_{?} \\
W \arrow[u] \arrow[r, "V"] & W \arrow[u] \arrow[r, no head] & W_1
\end{tikzcd}

\note{l'indice du dernier $W$ de la ligne du haut est illisible (rendu
$?$) ; les deux traits horizontaux de droite n'ont pas de pointe}

$w \in W_\infty$ \qquad $\mathbf{W}^{\xi}$
\note{devant le $\mathbf{W}$ gras, un signe noirci ; l'exposant est une
petite lettre ondulée, lue $\xi$}

\begin{quote}
$V(w)$

$V(wx) = \sigma^{-1}(w)\, V(x)$
\end{quote}
\note{ces deux formules sont dans un cadre ouvert, tracé à gauche et en
bas ; sous le trait du bas, un petit carré « $W_3 \to W_4$ », avec
flèches verticales vers $W_{?}$ et $W_4$ marquées $\sigma^{-3}(w)$
et $\sigma^{-3}(w)$ (le second exposant est surchargé), et à côté
$\overline{\mathbb{G}_a}$}

\[ \overline{W}_1 \qquad \underline{\mathcal{U}} \subset V \]
\[ w_1 \quad \sigma^{-1}(w_2) \quad \sigma^{-2}(w_2) \quad \sigma^{-3}(w_3) \]
\[ \overline{W}_1 \subset \overline{W}_2 \subset \overline{W}_3 \subset
\overline{W}_4 \subset \cdots \qquad \overleftarrow{R} \]
\note{cette colonne est séparée du reste par un double trait oblique ;
les indices $2$, $2$, $3$ de la deuxième ligne sont surchargés. La
flèche $\overleftarrow{R}$ est tracée sous la chaîne d'inclusions,
dirigée vers la gauche}

\struck{$F\,R = RF = p$}
\begin{quote}
$FV = VF = p$

$F^n \mathcal{U} = 0$

$p^n \mathcal{U} = 0$
\end{quote}
\note{ces trois relations sont encadrées d'un parallélogramme au crayon
au bas à droite ; la ligne biffée est au-dessus du cadre. Le
$\mathcal{U}$ est une grande lettre cursive $U$ ; le « $V$ » de
« $\mathcal{U} \subset V$ » est gras}

\page{199}

\[ A = W_\infty[[F, V]] \qquad
\left\lbrace
\begin{aligned}
FV &= VF = p \\
Fw &= \sigma(w)\, F \\
wV &= V \sigma(w)
\end{aligned}
\right. \]
\note{en haut à droite, isolé du calcul par un grand trait courbe. Dans
la troisième relation, un premier second membre est biffé avant
« $V\sigma(w)$ » ; la parenthèse fermante de $\sigma(w)$ n'est pas
ouverte sur la page}

\struck{$w + \sigma(w) F + \Sigma$}
\[ w + \sum_{1}^{\infty} w'_i F^i + \sum_{1}^{\infty} w''_j V^j \]
\[ \sum_{0}^{\infty} \sigma(w'_i) F^{i+1} +
\sum_{0}^{\infty} \sigma(w'_i) F^{i+1} V \]
\[ \sum_{0}^{\infty} \sigma^{-1}(w''_j)\, V F^i +
\sum_{0}^{\infty} \sigma^{-1}(w''_j)\, V F^i V \]
\[ w + \sum_{1}^{\infty} w'_i F^i + \sum_{1}^{\infty} w''_j V^j \]
\[
\begin{aligned}
&p\sigma(w''_1) + \Bigl(\sigma(w) F + \sum_{2}^{\infty} \sigma(w'_{i-1}) F^{i}\Bigr)
+ \Bigl(\sum_{1}^{\infty} p\,\sigma(w''_{j+1}) V^j\Bigr) \\
&p\sigma^{-1}(w'_1) + \Bigl(\sum_{1}^{\infty} p\,\sigma^{-1}(w'_{i+1}) F^i\Bigr)
+ \Bigl(\sigma^{-1}(w) V + \sum_{1}^{\infty} \sigma^{-1}(w''_j) V^{j+1}\Bigr)
\end{aligned}
\]
\note{calcul des produits à gauche et à droite par $F$ et par $V$ d'un
élément de $A$. Plusieurs indices sont surchargés et de lecture
incertaine : l'indice de $w''_1$ au début de la première des deux
dernières lignes, l'exposant de $F$ dans la première somme (un $i$
corrigé), la borne inférieure de la somme en $V^j$ (surchargée en $1$) ;
les deux premières lignes du calcul portent des bornes $0$ ou $1$
surchargées. Dans la deuxième ligne, la seconde somme est biffée d'un
trait par l'auteur ; les indices $w'_i$, $w''_j$ des lignes 3 et 4 sont lus d'après
la ligne 5}

$C$ \qquad \struck{\ill{}} \emph{artiniens} \add{\emph{hyperalgèbres}}
\note{« hyperalgèbres » est écrit au-dessus de « artiniens » ; un mot
biffé précède. Un trait horizontal traverse ensuite la page}

\begin{quote}
pro(\emph{groupes formels \add{\uncertain{artiniens}} sur $k$})

\emph{groupes formels artiniens sur $k$}
\end{quote}
\note{« artiniens » est écrit au-dessus de la première ligne, entre
parenthèses}

\begin{quote}
groupes alg.\ commutatifs finis (d'hyperalgèbre $U$)
$\Longleftrightarrow$ Modules de longueur finie sur $A$
(\uncertain{isomorphisme} $\mathrm{Hom}(M, I)$)

\struck{\ill{}}(groupes alg.\ finis) (noeth.) $\Longleftrightarrow$
Modules loc.\ de longueur finie sur $A$ (art.)

$\Updownarrow$ \uncertain{par} \uncertain{dualité}

« groupes formels » (art.) $\Longleftrightarrow$ Modules
\struck{\ill{}} \struck{(noeth.)} (noeth.\ i.e.\ mod.\ de type fini sur $A$)
\end{quote}
\note{bas de la page, écrit en partie par-dessus l'en-tête imprimé
renversé de la feuille ; le signe biffé devant « (groupes alg.\ finis) »
est sans doute un « pro » raturé}

\page{200}

\begin{quote}
groupes formels raisonnables \emph{\uncertain{sans} \uncertain{isogénie}}

$\Updownarrow$ \uncertain{contrav.}

pro-groupes alg.\ finis raisonnables \emph{sans}
\struck{\uncertain{isogénie}} \emph{\uncertain{alg.} finie}
\quad $\Longleftrightarrow$ \uncertain{cov.}\quad \emph{Modules de torsion}
\struck{\ill{}} \emph{de type fini} $A_{\mathbf{F}}$
$=$ \emph{modules de type fini sur $A_{\mathbf{F}}$}
\end{quote}
\note{une flèche courbe, marquée « \uncertain{anti} », descend du premier
énoncé vers la colonne de droite. Les lectures « sans isogénie »,
« contrav. » et « cov. » sont incertaines ; les deux lignes du bas à
droite sont écrites en colonne étroite, chaque mot souligné}

\[ \underline{\mathbf{W}} \to \underline{A}_{\mathbf{F}} = \mathcal{A}
\qquad W_{\sigma}[F, \mathbf{V}] \qquad \mathbf{W}_n \]
\note{un double trait oblique sépare $\mathcal{A}$ de
$W_{\sigma}[F, \mathbf{V}]$ et descend, par un trait brisé, vers
$\mathbf{W}_n$, précédé d'un signe raturé. La lettre après « $=$ »
est une grande lettre ronde, rendue $\mathcal{A}$}

$\mathbf{F}\ \mathbf{V}$

\begin{quote}
Modules de type fini $M$ sur $\mathbf{W}$, munis d'une application
\emph{\uncertain{additive} \uncertain{bijective}} $F : M \to M$, et d'une
application additive $V : M \to M$, telles que
\[ \left\lbrace
\begin{aligned}
F(wx) &= F(w)\, F(x) \\
V(wx) &= F^{-1}(w)\, V(x)
\end{aligned}
\right. \]
\end{quote}
\note{énoncé encadré. « additive bijective » est inséré au-dessus de la
ligne, avec un renvoi. Au-dessous du cadre, au crayon, une lettre $R$,
et un $\mathbf{W}$ gras}

\end{document}
